We present a set of simple theoretical results 
that demonstrate the optimality of treatment disparity, 
and highlight some properties of DLPs.  
% Our optimality results are all derived in the population 
% or \emph{infinite data} setting, 
% and assume knowledge 
% of the true conditional probability function $p_{Y|X, Z}(\vct x,z) \equiv \P(Y = 1 \mid X = \vct x, Z = z)$. 
% The main results of this section can be summarized as follows:
We summarize our results as follows:
% \subsection{Optimality of disparate treatment under oracle model}
\begin{enumerate}
\item Direct treatment disparity on the basis of $z$ 
is the optimal strategy 
for maximizing classification accuracy\footnote{Our results are all presented in terms of a more general performance metric, of which classification accuracy is a special case.}
subject to CV and $p$-\% constraints.  
\item When $X$ fully encodes $Z$, a sufficiently powerful DLP is equivalent to treatment disparity.
\end{enumerate}
In Section \ref{section:empirical}, 
we empirically demonstrate a related point:
\begin{enumerate}[resume]
 \item When $X$ only partially encodes $Z$,
a DLP may be suboptimal and can induce intra-group disparity
on the basis of otherwise irrelevant features correlated with $Z$.
\end{enumerate}


% \vspace{5px}
\paragraph{Treatment disparity is optimal}
Absent impact parity constraints, the Bayes-optimal decision rule for minimizing expected $0-1$ loss (i.e., maximizing accuracy) is given by 
%
%  \[
%  d_{\mathrm{uncon}}^*(\vct x,z) = \begin{cases}
%  	1 & p_{Y|X, Z}(\vct x,z) \ge 0.5 \\
%     0 & \text{otherwise}
%  \end{cases}.  
%  \]
$d_{\mathrm{uncon}}^*(\vct x,z) = \delta( p_{Y|X, Z}(\vct x,z) \ge 0.5)$, where $\delta()$ is an indicator function.
% 
% Then to maximize accuracy, 
% without any consideration of disparate impact,
% the Bayes-optimal predictor
% outputs $\hat{y}_i = 1$ for each example $i$ such that $\hat{p}^*(x_i) > .5$, and outputs $\hat{y}_i=0$ otherwise.
% In this section,

We now show that the optimal decision rules 
in the CV and $p$-\% constrained problems 
have a similar form.  
The optimal decision rule 
will again be based on thresholding $p_{Y|X, Z}(\vct x,z)$, 
but at \textit{group-specific thresholds}. 
%
These rules can be thought of as operationalizing 
the following mechanism: 
Suppose that we start with the classifications 
of the unconstrained rule $d_{\mathrm{uncon}}^*(\vct x,z)$, 
and this results in a CV gap of $q_a - q_b  > \gamma$.  
To reduce the CV gap to $\gamma$ 
we have two mechanisms: We can
(i)  flip predictions from $0$ to $1$ in group $b$, and (ii) we can flip predictions from $1$ to $0$ in group $a$. 
The optimal strategy is to perform these flips on group $b$ cases that have the highest value of $p_{Y|X, Z}(\vct x,z)$ 
and group $a$ cases that have the lowest value of $p_{Y|X, Z}(\vct x,z)$.

The results in this section adapt the work of \cite{corbett2017algorithmic}, 
who establish optimal decision rules $d$ 
under 
% different kinds of 
% fairness-related 
% constraints.  
exact parity.
In that work, the authors characterize 
the optimal decision rule $d = d(\vct x, z)$ 
that maximizes the \textit{immediate utility} 
$u(d, c) = \E[Yd(X,Z) - cd(X,Z)]$ for $(0 < c < 1)$,
under different exact parity criteria.  
We begin with a lemma showing 
that expected classification accuracy 
has the functional form 
of an immediate utility function.

\begin{lemma} \label{lemma:accuracy}
Optimizing classification accuracy is equivalent 
to optimizing immediate utility with $c = 0.5$. 
\end{lemma}

\begin{proof}
The expected accuracy of a binary decision rule $d(X)$ 
can be written as $\E[Yd(X) + (1-Y)(1-d(X))]$.  
Expanding and rearranging this expression gives
% \small
$$
\E[Yd(X) + (1-Y)(1-d(X))] = \E(2Yd(X) - d(X)) + \E(Y) + 1
	= 2u(d, 0.5) + \E(Y) + 1.
$$
% \normalsize
%Since we are assuming that $Y$ is not affected by $d$, 
The only term in this expression that depends on $d$ 
is the immediate utility $u$.  
Thus the decision rule that maximizes $u$ 
also maximizes accuracy.
\end{proof}

% Before proceeding to the main technical results, 
We note that the results in this section 
are related to the recent independent work 
of \cite{pmlr-v81-menon18a},
% In their paper, the authors 
who derive Bayes-optimal decision rules 
under the same parity constraints we consider here, 
working instead with the \textit{cost-sensitive risk}, 
$
\mathrm{CS}(d;c) = \pi(1-c)\mathrm{FNR}(d) + (1-\pi)c\mathrm{FPR}(d),
$
where $\pi = \P(Y = 1)$.  
One can show that $u(d, c) = -\mathrm{CS}(d;c) + \pi(1-c)$, 
and hence the problem of maximizing immediate utility considered here 
is equivalent to minimizing cost-sensitive risk 
as in \cite{pmlr-v81-menon18a}.  
In our case, it will be more convenient 
to work with the immediate utility.  

For the next set of results,
we follow \cite{corbett2017algorithmic} 
and assume that $p_{Y\mid X,Z}(X,Z)$, 
viewed as a random variable, 
has positive density on $[0,1]$.  
This ensures that the optimal rules 
are unique and deterministic 
by disallowing point-masses of probability 
that would necessitate tie-breaking 
among observations with equal probability.  
The first result that we state 
is a direct corollary 
of two results in \cite{corbett2017algorithmic}.  
It considers the case 
where we desire exact parity, i.e., 
that $q_a = q_b$.  
%Similarly, define $p_{Y|X} = p_{Y|X}(x) = \E(Y = 1 \mid X = x)$.

\begin{cor}
 The optimal decision rules $d^*$ 
 under various parity constraints 
 have the following form 
 and are unique up to a set of probability zero: 
 \begin{enumerate}
 \item Among rules satisfying statistical parity (the 100\% rule), the optimum is
%  \[
%  d^*(\vct x,z) = \begin{cases}
%  	1 & p_{Y|X, Z}(\vct x,z) \ge t_z \\
%     0 & \text{otherwise}
%  \end{cases},
%  \]
$d^*(\vct x,z) = \delta(p_{Y|X, Z}(\vct x,z) \ge t_z),$
 where $t_z \in [0,1]$ are constants 
 that depend only on group membership $z$. 
 \item Among rules that have equal false positive rates across groups, the optimum is
%  \[
%  d^*(\vct x,z) = \begin{cases}
%  	1 & p_{Y|X, Z}(\vct x,z) \ge s_z \\
%     0 & \text{otherwise}
%  \end{cases},
%  \]
 $d^*(\vct x,z) = \delta(p_{Y|X, Z}(\vct x,z) \ge s_z),$
 where $s_z$ are constants 
 that depend only on group membership $z$ 
 (but are different from $t_z$).
 \item (1) and (2) continue to hold 
 even in the resource-constrained setting 
 where the overall proportion of cases classified as positive is constrained.
 \end{enumerate}
\end{cor}

\begin{proof}
(1) and (2) are direct corollaries of Lemma \ref{lemma:accuracy} combined with Thm 3.2 and Prop 3.3 of \cite{corbett2017algorithmic}. 
\end{proof}

The next set of results establishes optimality 
under general $p$-\% and CV rules.

\begin{prop}
	Under the same assumptions as above, 
    the optimum among rules 
    that satisfy the CV constraint $0 \le q_a - q_b < \gamma$ 
    or the $p$-\% rule also has the form
%      \[
%  d^*(\vct x,z) = \begin{cases}
%  	1 & p_{Y|X, Z}(\vct x,z) \ge t_{z} \\
%     0 & \text{otherwise}
%  \end{cases},
%  \]
$d^*(\vct x,z) = \delta(p_{Y|X, Z}(\vct x,z) \ge t_{z}),$
 where $t_{z} \in [0,1]$ are constants 
 that depend on the group membership $z$, 
 and on the choice of constraint parameter $\gamma$ or $p$.  
 The thresholds $t_z$ are different 
 for the CV constraint and $p$-\% rule.
\label{prop:cv_opt}
\end{prop}

\begin{proof}
Suppose that the optimal solution under the CV or $p$-\% rule constraint classifies proportions $q_a$ and $q_b$
of the advantaged and disadvantaged groups, respectively,
to the positive class. 
As shown in \citet{corbett2017algorithmic}, 
we can rewrite the immediate utility as 
\[
 u(d,0.5) = \E[d(X, Z)(p_{Y|X, Z} - 0.5)].
\]
Thus the utility will be maximized 
 when $d^*(X,Z) = 1$ 
for the $q_z$ proportion of individuals in each group 
that have the highest values of $p_{Y|X, Z}$.  
Since the optimal values of $q_z$ 
may differ between the CV-constrained solution 
and the $p$-\% solution, 
the optimal thresholds may differ as well.
\end{proof}

%The final result in this section 
Our final result
shows that a decision rule 
that does not directly use $z$ as an input variable 
or for determining thresholds 
will have lower accuracy 
than the optimal rule 
that uses this information.  
That is, we show that DLPs are suboptimal 
for trading off accuracy and impact parity.

\begin{theorem}
	Let $d^*(\vct x, z)$ be the optimal decision rule under a the CV-$\gamma$ or $p$-\% constraint. 
    Let $d_{\mathit{DLP}}(\vct x)$ be the optimal solution to a DLP. 
        If $d(\vct x,z)$ and $d_{\mathit{DLP}}(\vct x)$ satisfy CV or $p$-\% constraints 
        with the same $q_a$ and $q_b$, 
        the DLP solution results in lower or equal accuracy 
        (equal only if the solutions are the same.)
%     \end{enumerate}
\end{theorem}

\begin{proof}
 From Proposition \ref{prop:cv_opt},
 we know that the unique accuracy-optimizing solution is given by 
%  \[
%  d^*(\vct x,z) = \begin{cases}
%  	1 & p_{Y|X, Z}(\vct x,z) \ge t_{z} \\
%     0 & \text{otherwise}
%  \end{cases},
%  \]
 $d^*(\vct x,z) = \delta(p_{Y|X, Z}(\vct x,z) \ge t_{z}),$
 where $t_{z}$ is the 1 - $q_{z}$ quantile of $p_{Y|X, Z}$.  
 The difference in immediate utility 
 between the two decision rules 
 can be expressed as follows:
 \small
 \begin{align*}
 	&\E[d^*(X, Z)(p_{Y|X, Z} - .5)] - \E[d_{\mathit{DLP}}(X)(p_{Y|X, Z} - .5)] \\
     &= \E[(d^*(X, Z) - d_{\mathit{DLP}}(X))(p_{Y|X, Z} - 0.5)] \\
     &= \E[p_{Y|X, Z}\!-\! \!.5\!\! \mid\!\! d^*\! =\! 1, d_{\mathit{DLP}}\! =\! 0]\P(d^*\!\! =\! 1, d_{\mathit{DLP}}\! =\! 0)
     \!-\!\E[p_{Y|X, Z}\! -\!\! .5\! \!\mid\!\! d^*\!\! =\! 0, d_{\mathit{DLP}}\! =\! 1]\P(d^*\!\! =\! 0, d_{\mathit{DLP}}\! =\! 1) \\
     &= \bigl( \E[p_{Y|X, Z} - .5 \mid d^* = 1, d_{\mathit{DLP}} = 0]  - \E[p_{Y|X, Z} - .5 \mid d^* = 0, d_{\mathit{DLP}} = 1] \bigr) \P(d^* = 1, d_{\mathit{DLP}} = 0) \\
     &\ge 0
 \end{align*}
 \normalsize
 The final inequality follows 
 %from the observation 
 %that 
 since
 $d^*(X, Z) = 1$ for the highest values of $p_{Y|X, Z}$, 
 so $p_{Y|X,Z}$ is stochastically greater 
 on the event $\{d^*\!=\! 1, d_{\mathit{DLP}}\! =\! 0\}$ 
 than on $\{d^*\! =\! 0, d_{\mathit{DLP}}\! =\! 1\}$.  
 Note that equality holds only if 
 $\P(d^*\! =\! 1, d_{\mathit{DLP}}\! =\! 0) = 0$, i.e., 
 if the two rules are (almost surely) equivalent.
\end{proof}

%All results in this section 
Our results
continue to hold 
under ``do no harm'' constraints, where we require that any individual in the disadvantaged group who was classified as positive under 
% where the proportion of cases 
% in the disadvantaged group classified as positive 
% is constrained to be no lower 
% than the proportion under 
the unconstrained rule $d_\mathrm{uncons}(\vct x,z)$ remains positively classified.  This corresponds to the setting where 
the proportion of cases 
in the disadvantaged group classified as positive 
is constrained to be no lower 
than the proportion under the unconstrained rule 
%$d_\mathrm{uncons}(\vct x,z)$
(or no lower than some fixed value $q_a^\mathrm{min}$).  
Such constraints impose an upper bound on the optimal thresholds $t_b$, but do not change the structure of the optimal rules. 


% \vspace{5px}
\paragraph{Functional equivalence when protected characteristic is redundantly encoded.}
% \paragraph{Functional equivalence when protected characteristic is redundantly encoded} 
Consider the case where the protected feature $z$ 
is redundantly encoded in the other features $\vct x$.  
More precisely, 
suppose that there exists a known subcomputation $g$ 
such that $z = g(\vct x)$.  
This allows for any function of the data $f(\vct x,z)$ 
to be represented as a function of $\vct x$ alone via $\tilde f(\vct x) = f(\vct x, g(\vct x))$.   
While it remains the case that $\tilde f(\vct x)$ does not directly use $z$ as an input variable---and thus satisfies treatment parity---$\tilde f$ should be no less legally suspect from a \textit{disparate treatment} perspective 
than the original function $f$ that uses $z$ directly. 
The main difference for the purpose of our discussion 
is that $\tilde f$, resulting from a DLP, may technically satisfy treatment parity, while $f$ does not.  
%Yet it is clear that $f$ and $\tilde f$ are functionally equivalent in two senses: (i) as mathematical objects; and more importantly, (ii) as mechanisms for classifying cases. 

While this form of ``strict'' redundancy is unlikely, characterizing this edge case is
important for considering whether DLPs should have different legal standing vis-a-vis disparate treatment than methods
that use $z$ directly. This is particularly relevant if one thinks of the `practitioner' in question as having discriminatory
intent. Furthermore, the partial encoding of the protected attribute is commonplace in settings where discrimination
is a concern (as with gender in our experiment in \S \ref{section:empirical}). Indeed, the very premise of DLPs requires that $\vct x$ is
significantly correlated with $z$. Moreover, DLPs provide an incentive for practitioners to game the system by adding
features that are predictive of the protected attribute but not necessarily of the outcome, as these would improve the
 DLP's performance.

\paragraph{Within-class discrimination when protected characteristic is partially redundantly encoded.}
% \paragraph{Within-class differentiation when protected characteristic is partially redundantly encoded} \label{subsec:partial}
% \note{Add a simple analytic example where there are two features, one of which is uncorrelated with $z$ and the other is uncorrelated with $y$.}
When the protected characteristic is partially encoded in the other features, disparate treatment may induce within-class discrimination by applying the benefit of the affirmative action unevenly, and can even harm some members of the protected class.
%This claim asserts a possibility, so it is sufficient to produce one example to support the claim. 
%In the following section, 
Next
we demonstrate this phenomenon empirically using 
(synthetically biased) university admissions data and several public datasets. 
The ease of producing such examples might convince the reader that the 
%highly 
varied effects of intervention with a DLP 
on members of the disadvantaged group raises practice and policy concerns about DLPs.

% \zack{Argument sketch}
% \begin{itemize}
% \item show that optimal thing to do is to 
% estimate $P(y)$ and $P(z)$ separately. \item Given $P(y)$ and $P(z)$ optimal decisions can be made with a trivial extension to the scoring mechanism from subsection 1. 
% \item Implication - a powerful enough model will implicitly learn $P(z)$.
% \item Can we demonstrate how exactly this might happen, even with say a linear model?
% \end{itemize}
% We demonstrate this phenomenon empirically the next section.
